\begin{proof}
We establish joint stability, give exact error Gramians, and differentiate
the isolated leading squared singular value.

\emph{Joint stability.}
Direct multiplication yields
\begin{equation}
 A_a^*A_a+A_b^*A_b
   =\operatorname{diag}(1-q-q^2,1).
 \label{I:F:eq:f2-passivity}
\end{equation}
Let $D_0=\operatorname{diag}(q+q^2,0)$. The matrix
$D_0+A_a^*D_0A_a+A_b^*D_0A_b$ is positive definite. Indeed, its kernel
would be contained in $\ker D_0=\operatorname{span}\{(0,1)^{\mathsf T}\}$, but
$A_b(0,1)^{\mathsf T}=(1,1)^{\mathsf T}/\sqrt2$ has nonzero first coordinate.
Telescoping \eqref{I:F:eq:f2-passivity} therefore gives
\[
 \sum_{|w|=2}A_w^*A_w\prec I.
\]
The column completely positive map is a strict contraction after two
steps, and its adjoint has the same spectral radius. Thus the system is
mean-square stable. Its word-Hankel is bounded of rank at most two.

\emph{Competitor Gramians.}
Let $R_g,Q_g$ be its exact reachability and observation Gramians. Their
Stein equations have the solutions
\begin{equation}
 R_g=\begin{pmatrix}
 \dfrac{1+q^2}{q(1+q)}&\dfrac{1-q}{1+q}\\[4pt]
 \dfrac{1-q}{1+q}&\dfrac{2+q+2q^2+q^3}{1+q}
 \end{pmatrix},\qquad
 Q_g=\frac{q(1+2q)^2}{4(1+q)}I_2.
 \label{I:F:eq:f2-own-gramians}
\end{equation}
These identities can be verified by substitution in
$R_g=\beta_g\beta_g^*+\sum_j A_jR_gA_j^*$ and
$Q_g=\alpha_g\alpha_g^*+\sum_j A_j^*Q_gA_j$.
For example, writing the entries of $R_g$ as $u,v,w$, the off-diagonal
and lower-right equations give
$v=(1-q)/(1+q)$ and $w=2+q^2u-2qv$; the remaining equation gives
$u=(1+q^2)/(q(1+q))$. Joint stability ensures uniqueness.

\emph{Target and cross Gramians.}
Use the target realization
\[
 F_a=\operatorname{diag}(r,0,0),\qquad
 F_b=\begin{pmatrix}0&1&0\\0&0&1\\0&0&0\end{pmatrix},\qquad
 \alpha_f^*=q(1,0,0),\qquad \beta_f=(0,0,1)^{\mathsf T}.
\]
Its Gramians are
$R_f=\operatorname{diag}(q^{-1},1,1)$ and $Q_f=qI_3$.
Put $\Delta=1-ra>0$. Define $X,Y\in\FC^{3\times2}$ by their rows:
\begin{align}
 X_0&=\left(\frac{1+q+q^2-q^3}{2\Delta},
                 \frac{1-q+q^2+q^3}{2}\right),\notag\\
 X_1&=\frac1{\sqrt2}(1-q^2,1+q^2),\qquad X_2=(-q,1),
 \label{I:F:eq:f2-cross-reach}\\
 Y_0&=\frac{qc}{\Delta}(1,0),\qquad
 Y_1=Y_0A_b,\qquad Y_2=Y_0A_b^2.
 \label{I:F:eq:f2-cross-observe}
\end{align}
The triangular form of $F_b$ gives directly
\[
 X=\beta_f\beta_g^*+\sum_jF_jXA_j^*,\qquad
 Y=\alpha_f\alpha_g^*+\sum_jF_j^*YA_j.
\]
For example, the first equation gives successively
$X_2=\beta_g^*$, $X_1=\beta_g^*A_b^*$, and
$X_0=\beta_g^*(A_b^*)^2\operatorname{diag}(\Delta^{-1},1)$,
which is \eqref{I:F:eq:f2-cross-reach}.

For the composite error realization with output
$(\alpha_f^*,-\alpha_g^*)$ and input $(\beta_f,\beta_g)$, the exact
Gramians are therefore
\begin{equation}
 R=\begin{pmatrix}R_f&X\\X^*&R_g\end{pmatrix},\qquad
 Q=\begin{pmatrix}Q_f&-Y\\-Y^*&Q_g\end{pmatrix}.
 \label{I:F:eq:f2-error-gramians}
\end{equation}
Its infinite Hankel squared norm is the largest eigenvalue of
$R^{1/2}QR^{1/2}$, equivalently the largest eigenvalue of $RQ$.
Equations \eqref{I:F:eq:f2-own-gramians}--\eqref{I:F:eq:f2-error-gramians}
are thus a finite exact algebraic description of this chosen error;
there is no word-horizon cutoff.

\emph{Leading eigenvalue perturbation.}
Set $\widehat R=qR$ and $\widehat Q=Q/q$, so
$\widehat R\widehat Q=RQ$. Both extend analytically to $q=0$.
In particular $\Delta=q+q^2+O(q^4)$, removing the apparent
singularities in $qX$ and $Y/q$.
Index the five coordinates by $(f_0,f_1,f_2,g_0,g_1)$ and put
$y=1/\sqrt2$. The expansions are
\[
 \widehat R=R_0+qR_1+O(q^2),\qquad
 \widehat Q=Q_0+qQ_1+O(q^2),
\]
where
\begin{equation}
 R_0=\begin{pmatrix}
 1&0&0&1/2&0\\0&0&0&0&0\\0&0&0&0&0\\
 1/2&0&0&1&0\\0&0&0&0&0
 \end{pmatrix},\quad
 R_1=\begin{pmatrix}
 0&0&0&0&1/2\\0&1&0&y&y\\0&0&1&0&1\\
 0&y&0&-1&1\\1/2&y&1&1&2
 \end{pmatrix},
 \label{I:F:eq:f2-R-expansion}
\end{equation}
\begin{equation}
 Q_0=\begin{pmatrix}
 1&0&0&-1/2&0\\0&1&0&0&-y/2\\0&0&1&0&-1/4\\
 -1/2&0&0&1/4&0\\0&-y/2&-1/4&0&1/4
 \end{pmatrix},\quad
 Q_1=\begin{pmatrix}
 0&0&0&-1/2&0\\0&0&0&-y/2&-y/2\\0&0&0&1/4&-1/2\\
 -1/2&-y/2&1/4&3/4&0\\0&-y/2&-1/2&0&3/4
 \end{pmatrix}.
 \label{I:F:eq:f2-Q-expansion}
\end{equation}
The matrix $Z_0=R_0Q_0$ has only one nonzero row:
$Z_0[0,0]=3/4$, $Z_0[0,3]=-3/8$.
Its sole nonzero eigenvalue $3/4$ is simple. Right and left eigenvectors,
normalized by $w^*v=1$, are
\[
 v=e_0,\qquad w=e_0-\tfrac12 e_3.
\]
The analytic eigenvalue branch consequently has derivative
\begin{equation}
 \lambda'(0)=w^*(R_1Q_0+R_0Q_1)v=-\frac14.
 \label{I:F:eq:f2-eigen-derivative}
\end{equation}
More explicitly, the two terms are $-1/4$ and zero, respectively.
All other eigenvalues tend to zero, so this branch is the largest
squared Hankel singular value for all sufficiently small positive $q$.
Analyticity supplies a bounded $O(q^2)$ remainder and proves
\eqref{I:F:eq:f2-branch-expansion}. Admissibility of the real two-state
construction proves \eqref{I:F:eq:f2-optimum-upper}.
\end{proof}
